\chapter{まとめと今後の方針}

本章では，現時点までに得られた成果をまとめ，卒業論文の完成に向けた今後の方針を述べる．

\section{現時点のまとめ}

本論文の時点で得られた成果は 3 点である．

第一に，3 つの確信度引き出し方式を設計し，これらを同一条件で比較する評価基盤を実装した．基盤は ECE，Brier Score，AUROC，Murphy 分解を計算し，問題の復元抽出による信頼区間を算出する．応答全文を保存する設計としたため，採点処理の不具合を追加の問い合わせなしに修正できた．

第二に，日本語版 MGSM の 250 問を用いた予備実験を行った．Ling.1S が他の 2 方式より ECE，Brier Score，AUROC，正答率で劣ることが信頼区間によって支持された一方，REL と RES では差を検出できず，Verb.1S と Verb.2S の差はどの指標でも判定できなかった．

第三に，3 方式に共通して確信度が 0.9 以上に偏ることを観察した．この分布のもとでは ECE の値が平均確信度と平均正答率の差に近づき，較正の良否と正答率の高低を読み分けにくい．当初は Murphy 分解の REL によって読み分けが可能と考えていたが，REL も各区間の正答率を含むため分離の手段にはならなかった．

\section{今後の計画}

今後の作業の順序を表 6.1 に示す．

\begin{table}[htbp]
\centering
\small
\caption{今後の作業計画}
\label{tab:6.1}
\begin{tabularx}{\linewidth}{|l|X|l|}
\hline
順序 & 作業 & 追加の実行 \\
\hline
1 & 採点処理の検証と保存仕様の改善 & 不要 \\
2 & 言語表現の対応づけの感度確認 & 不要 \\
3 & 較正のみを取り出す比較条件の実行 & 要 \\
4 & 確信度を尋ねない基準条件の実行 & 要 \\
5 & プロンプト条件の設計と少数条件での検証 & 要 \\
6 & データセットと対象モデルの拡張 & 要 \\
\hline
\end{tabularx}
\end{table}

順序 1 と 2 を先に置くのは，追加の問い合わせを必要とせず，かつ以降のすべての結果の信頼性に関わるためである．採点が正しいことを確かめないまま条件を増やすと，誤りを含んだまま規模だけが大きくなる．応答の形式を分類して抽出と判定が意図通りかを確認し，Verb.2S の第 1 ターンの応答も保存するよう実装を改める．

順序 3 と 4 は，5.4 節で述べた較正と正答率を分離できないという問題に対処する．すべての方式に共通の回答をあらかじめ与えて確信度のみを尋ねれば，正誤ラベルが方式間で共通になるため較正の違いだけを比較できる．あわせて，確信度を尋ねずに回答のみを求める基準条件も設ける．

順序 5 が本研究の中心課題である．検討している方向は 3 つある．1 つ目は，確信度を答える前に，その回答が誤っている可能性のある理由を挙げさせる方向である．自信過剰が，誤りの可能性を検討しないまま確信度を述べることに由来するのであれば，検討を促すことで確信度が変わると期待される．2 つ目は，各段階を使う頻度に関する指示をプロンプトに加える方向であるが，較正を改善したのではなく分布を操作しただけになる恐れがある．3 つ目は，複数の問題を提示して確信の高い順に並べさせる方向であるが，順位を確率に変換する規則が要る．いずれの場合も，分布を広げるだけでは較正は改善しないため，偏りの解消そのものは改善の条件としない．

順序 6 では JCommonsenseQA と JMMLU による追試，および対象モデルの拡大を行う．順序 5 と 6 では，選択の影響が評価に入り込むことを避けるため，条件を選ぶためのデータと最終評価に用いるデータを分ける．
